\justifying
\NSFCBodyText

\subsubsection{项目特色}
项目以可产生环境状态变化的多步骤任务为研究对象，将上下文管理与行动所依赖的证据连接起来。静态长文理解和会话记忆作为诊断层，程序化终态测试作为主要验证层；同一轨迹同时记录状态、来源、版本与成本，使压缩效果能够追溯到具体的遗漏或恢复事件。方案采用冻结执行模型与渐进开发，在已有算力及一年期草案内优先完成能够被独立复核的机制实验。

\subsubsection{创新点}
\textbf{（1）研究不完备元数据下依赖保留与主动取证的预算配置。}现有轨迹压缩、原文回读及版本化记忆已提供重要基础。已有研究还覆盖记忆失效、动作纠错和回放。本项目聚焦依赖及版本信息不完全可见的条件，以未完成动作的前置条件连接压缩与恢复，分层改变压缩强度、依赖可见程度及元数据错漏，检验有限预算在状态保留与主动重查之间如何分配。以相同候选证据、摘要长度与恢复开销排除资源差异，使观察到的改善能够对应到具体任务条件；对静态、弱依赖和元数据不完备情境同时检验失效边界。创新成立与否取决于这一可证伪关系及针对性方法是否经验证，不以模块组合名称替代证据。

\textbf{（2）建立由恢复风险决定、覆盖全部控制开销的调度收益判据。}在已有预算上下文管理研究基础上，把前提缺失、版本不确定性和实际恢复代价共同用于决策，研究额外控制何时能带来净收益。通过任务级非劣比较与完整成本核算，区分降低窗口峰值、减少累计Token和提高任务可靠性三类效果；通过规则、轻量策略及不可部署理想参照揭示可改进空间。预期形成不同依赖跨度、噪声和更新频率下的适用依据，负结果亦作为明确方法边界的证据。
